% !TeX root = ../kazymyr-thesis.tex

\chapter{Conclusion and Future Work}\label{chapter:conclusion}

\begin{comment}
\todo{summary --- to be written once the remaining work is finished; the
  RQ1--RQ3 answers below are drafted from Chapters~\ref{chapter:approach}
  and~\ref{chapter:experiments} and are kept commented out until then}

% ---------------------------------------------------------------------
% Summary, hidden for now. Restore by uncommenting the block below.
% ---------------------------------------------------------------------
% \section{Summary}\label{sec:concl-summary}
%
% \subsection{RQ1: making the hybrid mode well defined}\label{sec:concl-rq1}
%
% ACID's lock table and BASE's dirty set are independent mechanisms that
% never coexist in a pure run. Selecting the mode per transaction allows a
% resource to be locked and dirty at the same time, which is contradictory.
% The shared-resource arbiter resolves this by making ACID wait on dirty
% while BASE defers its repair on locked. Because the only cross-mode wait edge points one
% way, no cross-mode wait cycle can form, and the arbiter introduces no
% deadlock that pure ACID did not already admit. The resolution is not free;
% its cost is quantified under RQ2.
%
% \subsection{RQ2: performance and tunability}\label{sec:concl-rq2}
%
% The hybrid mode offers a configurable inconsistency budget and charges for
% it.
%
% The adaptive oracle engages only inside an interior band of reconciler
% provisioning, $1 \le R \le 15$ at the baseline configuration, and is inert
% outside it; at $R = 0$ and $R \ge 25$ every statistic is identical to pure
% BASE in every seed. Inside that band, throughput does not interpolate
% between the two pure modes but falls below a time-weighted mixture of them
% by 2.8 to 12.3 times, with the gap widening as reconciliation slows.
%
% The critical dirty fraction $c$ nonetheless behaves as a configurable
% bound on the tail. Across $c$ from 2\% to 30\% and $R$ from 1 to 15,
% measured $p_{95}$ never departs from configured $c$ by more than
% 0.74 percentage points, with a median seed-to-seed coefficient of variation
% of 0.25\%. X4 and X5 extend the same result across two workload axes: it
% spans 1.04\,pp across semantic overlap from 0.20 to 0.50 and
% 0.78\,pp as the resources per transaction grow 20 times, in
% both
% cases straddling the configured value. It stops applying in two predictable and detectable ways: once
% $c$ exceeds the BASE steady-state dirtiness at that $R$, and once
% reconciliation capacity falls below the dirtying rate the workload
% imposes, at which point the oracle saturates in ACID and misses the budget
% regardless of configuration.
%
% \subsection{RQ3: is the feedback necessary?}\label{sec:concl-rq3}
%
% The answer depends on what is being asked for, and the two halves point
% in opposite directions.
%
% On the trade-off itself, feedback is not necessary and it is not free.
% Compared at the same measured inconsistency, the two oracles are level at
% $R = 1$ and $R = 5$, and the adaptive one is between a tenth and a fifth
% slower at $R = 10$ and $R = 15$ (Section~\ref{sec:x3-frontier}). The cause
% is measured rather than guessed. The adaptive oracle enters ACID at the
% moment the dirty fraction has just crossed $c$, so it admits its expensive
% transactions when the dirty set is at its largest, and a higher share of
% them block on a resource waiting to be repaired. Where repair is the limit
% that costs nothing, and where capacity is the limit it costs throughput.
%
% That result should be read with the workload in mind. Every run in this
% thesis is executed under a stationary load, and an adaptive controller has
% nothing to adapt to when the operating point never moves; the experiment
% was always going to be a poor place for feedback to show its value.
% Running the same comparison under a workload that changes while the run is
% in progress is the single most valuable extension of this work
% (Section~\ref{sec:future-work}), and until it is done, RQ3 is answered
% only for the stationary case.
%
% The other half is not a throughput result at all. Each oracle holds fixed
% the quantity it reads and lets the other float. The fixed-share oracle
% realises its configured ACID share to within 6\% across the whole
% reconciler range while its measured $p_{95}$ moves by up to 38.8 times,
% and the adaptive threshold is the mirror image, holding $p_{95}$ within
% 1.5 percentage points of the configured $c$ while its realised ACID share
% varies by 2.8 times. Only one of those two quantities is the one an
% operator wants to bound. A fixed share reaches a 10\% tail too, but the
% setting that does it depends on the provisioning and has to be found by
% measurement for each one; $c = 0.10$ holds it everywhere.
%
% The adaptive oracle therefore does not make the ACID/BASE trade-off
% cheaper. It makes it addressable, and at high reconciler counts that
% costs some throughput.
\end{comment}

\section{Lessons learned}\label{sec:concl-lessons}

This section collects the methodological lessons of the work. They are
about how the results had to be produced and read, rather than about the
hybrid mode itself, and each of them cost something before it was
understood.

\begin{itemize}
\item A HYBRID run with zero mode switches \emph{is} pure BASE with a
different label. Engagement must be verified before a result is
interpreted: five of the 24 cells in X2 are degenerate in exactly this
way, and thirteen of the 28 in X4.

\item A configured threshold bounds the estimate, not the underlying
quantity. The two coincide only when the system is slow relative to the
sampler, and when they diverge, the direction of the error depends on
which mode the condition is evaluated in.

\item Composing two correct mechanisms does not yield a mechanism whose
cost is the weighted average of theirs.

\item A metric with three plausible definitions needs all three named
before any of them is reported. Choosing one and describing another
produced a headline number that was wrong by six times, and the
error survived several readings because both quantities are fractions and
both were plausible in size.

\item An analytical model can be qualitatively right and mechanistically
wrong. The blocking probability derived in
Section~\ref{sec:approach-blocking-cost} is close to one across most of
the operating range, which the measurements confirm --- and that is
precisely why it explains nothing: a saturated variable has no variation
left with which to explain a response.

\item Comparing a controller against a baseline without feedback requires
deciding which variable is held constant. If that variable is an input to
one arm and an output of the other, the comparison is not well posed,
however carefully it is measured. The way out is to match on a quantity
that is an output of both arms and to interpolate the baseline onto it,
rather than to pair up whichever settings happen to land near each other:
matching to the nearest sampled cell discarded most of the data and left a
residual mismatch as large as the effect being measured.
\end{itemize}

\section{Future work}\label{sec:future-work}

This section lists the extensions the thesis leaves open, ordered roughly by
how much they would add. The first ones address limits of the evaluation,
the later ones address limits of the design itself.

\begin{enumerate}
\item \textbf{Non-stationary workloads} (terminal-based arrivals,
phase-changing overlap). This is the setting in which adaptive control
would be expected to earn back the throughput it costs under stationary
load (Section~\ref{sec:x3-frontier}), and it is the most valuable single
extension.

\item \textbf{A finer grid of ACID shares.} The fixed-share arm is
sampled at five values of $p_{\mathrm{acid}}$, and a third of the X3 comparisons fall too
far from any of them to be read on their own
(Section~\ref{sec:x3-frontier}). Sampling between $p_{\mathrm{acid}} = 0.10$ and
$p_{\mathrm{acid}} = 0.25$, and below $p_{\mathrm{acid}} = 0.05$ where the trade-off is steepest, would
turn those readings into measurements.

\item \textbf{Tick-driven signal sampling}, to decouple the filter's time
constant from the load it measures. Because admissions track completions
under \texttt{EQUILIBRIUM}, the oracle's sampling rate is approximately
the throughput its own decisions suppress, and the measured interval
varies by 60 times across the sweep.

\item \textbf{Within-visit sampling analysis.} The undershoot of the
recovery threshold at exit is measured, but the mechanism that makes it
grow with $R$ is not established, and the aggregate sample counts rule out
the obvious explanation. The tick timelines record per-tick state and
would settle it.

\item \textbf{A time-based oracle clock.} The oracle samples once per
admission, so its sampling rate collapses with the throughput it causes.
Sampling on a fixed interval of simulated time would decouple the two and
is the single change most likely to close the remaining exit slack.

\item \textbf{Absolute dirty thresholds} alongside fractional ones, for
small systems.

\item \textbf{Predictive switching}, on the dirty \emph{trend} rather than
the level.

\item \textbf{Lock-granularity aliasing} in the check that defers a
BASE repair.

\item \textbf{Real-system validation} against the Vitruvius case-study
parameters under HYBRID.

\item \textbf{Transaction-size interaction.} The blocking mechanism makes
the maximum transaction size $L$ a hidden co-determinant of the mixing
cost. Sweeping $L$ at fixed $c$ would test the mechanism directly: since
block duration rather than incidence carries the cost, $L$ should move the
number of dirty resources waited on, and therefore the block duration,
roughly linearly.

\item \textbf{Closing the fixed-share frontier at small $p_{\mathrm{acid}}$.} The
interval $p_{\mathrm{acid}} \in (0, 0.05)$ is unsampled and contains the steepest part of
the trade-off.

\item \textbf{Reporting the arbiter invariant empirically.} The
arbiter's invariant and the deadlock-freedom argument
(Sections~\ref{sec:approach-arbiter-rule} and~\ref{sec:approach-deadlock})
are established by construction and on paper; no measurement of them is
reported. The evidence already exists in the run data and would need no new
runs. Both sides of the rule are counted per run as
\texttt{acid\_txn\_blocked\_on\_dirty\_count} and
\texttt{base\_txn\_yielded\_on\_locked\_count}, and both are exactly zero in
every pure ACID and pure BASE run and at the degenerate adaptive points
$R = 0$ and $R = 25$, while in the engaged band they number in the hundreds
and tens of thousands respectively. Reporting them would show that the
mixed state was reached and resolved constantly rather than merely never
observed, which is what makes ``the invariant was never violated'' a claim
with content. It would also confirm the degeneracy of $R = 0$ and $R = 25$
through a mechanism independent of the switch count. Two cautions apply to
the wording: the invariant is enforced by construction at the two write
paths, not checked by a runtime assertion --- the Java \texttt{assert}
statements in the codebase were disabled in every experiment run, as no
harness passes \texttt{-ea} --- so the claim to make is that the hazardous
state cannot be constructed, evidenced by how often the enforcing branch was
taken. The natural home is
Section~\ref{sec:approach-arbiter}, after the deadlock argument, rather
than the experiments chapter, since neither property varies with $R$ or
$c$.
\end{enumerate}
